\documentclass[tikz,border=8pt]{standalone}
\usepackage{iftex}
\ifLuaTeX
  \usepackage{fontspec}
  \usepackage{unicode-math}
  \setmainfont{Libertinus Serif}
  \setsansfont{Libertinus Sans}
  \setmathfont{Libertinus Math}
\else
  \usepackage[T1]{fontenc}
  \usepackage{libertinus}
  \usepackage{libertinust1math}
\fi
\usepackage{amsmath}
\ifLuaTeX\else\usepackage{amssymb}\fi
\usepackage{fontawesome5}
\usetikzlibrary{positioning,arrows.meta,fit,backgrounds,calc,decorations.pathreplacing,shadows.blur,shapes.misc}

% ── Nature / Science journal-inspired palette ─────────────────
% Primary: deep navy (pipeline steps, badges)
\definecolor{navyD}{HTML}{1B2A4A}    % dark navy
\definecolor{navyM}{HTML}{2C5282}    % medium navy
\definecolor{navyL}{HTML}{EBF0F7}    % very light navy wash
% Accent: warm terracotta/rust (novel contributions: FPG + loss)
\definecolor{rustD}{HTML}{9B4D2B}    % deep terracotta
\definecolor{rustM}{HTML}{C76B3F}    % medium rust
\definecolor{rustL}{HTML}{FDF2EC}    % warm cream wash
\definecolor{rustE}{HTML}{D4845A}    % rust for edges
% Neutral: warm grey (inputs, data)
\definecolor{neutD}{HTML}{4A4A4A}    % charcoal
\definecolor{neutM}{HTML}{6B7280}    % medium grey
\definecolor{neutL}{HTML}{F5F5F3}    % warm off-white
% Sage green (categorization layer)
\definecolor{sageD}{HTML}{3D6B5E}    % deep sage
\definecolor{sageM}{HTML}{5A9E8F}    % medium sage
\definecolor{sageL}{HTML}{EDF5F2}    % sage wash
% Text
\definecolor{dk}{HTML}{1A1A2E}       % near-black with warmth
\definecolor{gy}{HTML}{4A4A4A}       % dark grey for annotations
\definecolor{arwC}{HTML}{2C3E50}     % arrow color

\begin{document}
\begin{tikzpicture}[
  >=Stealth,
  every node/.style={font=\sffamily\small, text=dk},
  % ── Inputs (warm neutral) ──
  inpMain/.style={
    draw=neutD!60, fill=neutL, rounded corners=4pt,
    minimum height=0.85cm, font=\sffamily\normalsize, align=center,
    inner xsep=10pt, line width=0.7pt},
  inpSub/.style={
    draw=neutD!35, fill=neutL!80, rounded corners=3pt,
    minimum height=0.65cm, minimum width=4.2cm,
    font=\sffamily\footnotesize, align=center, inner xsep=8pt,
    line width=0.5pt},
  % ── Processing steps (deep navy) ──
  stpbox/.style={
    draw=navyM!80, fill=navyL, rounded corners=4pt,
    minimum height=0.9cm, minimum width=12.4cm,
    font=\sffamily\normalsize, align=center, inner xsep=8pt,
    line width=0.7pt},
  stlbl/.style={font=\sffamily\normalsize, text=dk},
  % ── Novel: FPG feature groups (neutral, matching inputs) ──
  grp/.style={
    draw=neutD!50, fill=neutL, rounded corners=2.5pt,
    minimum width=2.35cm, minimum height=0.6cm,
    font=\sffamily\small, align=center, inner sep=2pt,
    line width=0.6pt},
  % ── Fault category chips (navy, matching steps) ──
  catchip/.style={
    draw=navyM!60, fill=navyL, rounded corners=2.5pt,
    minimum width=2.85cm, minimum height=0.6cm,
    font=\sffamily\small, align=center, inner sep=2pt,
    line width=0.6pt},
  % ── Output (double border, navy wash — conclusion emphasis) ──
  outbox/.style={
    double=navyL!70, double distance=1.5pt, draw=navyM!60,
    rounded corners=4pt,
    font=\sffamily\normalsize, align=center,
    inner xsep=14pt, inner ysep=10pt, line width=0.6pt},
  % ── Loss function boxes ──
  lossbox/.style={
    fill=rustL!60, rounded corners=4pt, draw=rustD!35,
    font=\sffamily\small, align=center, inner xsep=12pt, inner ysep=10pt,
    text=dk, line width=0.5pt},
  % ── Loss labels (plain text, no box) ──
  losspill/.style={
    font=\sffamily\small, text=gy,
    inner xsep=2pt, inner ysep=2pt},
  % ── Healthy terminal (plain text, no box) ──
  healthy/.style={
    font=\sffamily\small, text=sageD,
    inner xsep=2pt, inner ysep=2pt},
  % ── Annotations & arrows ──
  annot/.style={font=\sffamily\small\itshape, text=gy},
  mainArrow/.style={->, thick, arwC!90},
  condArrow/.style={->, thick, arwC!90},
  badge/.style={
    circle, draw=navyM!80, fill=navyL,
    font=\sffamily\small\bfseries, text=navyM,
    minimum size=0.52cm, inner sep=0pt, line width=0.7pt},
  % ── FPG propagation edges (directed) ──
  fpgEdge/.style={-{Stealth[length=3.5pt, width=2.5pt]}, neutM, line width=0.7pt},
]

% ══════════════════════════════════════════════════════════════════
% TRANSFORMER MODEL
% ══════════════════════════════════════════════════════════════════
\node[inpMain, minimum width=5.4cm, font=\sffamily\normalsize]
  (transformer) at (0, 3.6) {{\textcolor{navyM}{\faMicrochip}}\enspace Transformer Model};
\node[font=\sffamily\small, text=gy] at (0, 2.9)
  {(encoder or decoder)};

% Short arrow down from subtitle
\draw[mainArrow] (0, 2.65) -- (0, 2.0);

% Two branches with proper icons
\node[inpSub, draw=neutD!20, fill=neutL!30, text=gy!60] (static)
  at (-3.2, 1.1) {{\textcolor{neutM!50}{\faCode}}\enspace Source Code (Static)};
\node[inpSub] (dynamic)
  at ( 3.2, 1.1) {{\textcolor{neutD}{\faChartLine}}\enspace {Training Run} (Dynamic)};

% Static branch: dashed (not used in current work)
\draw[->, thick, arwC!30, densely dashed] (0, 2.0) -| (static.north);

\draw[mainArrow] (0, 2.0) -| (dynamic.north);

% ══════════════════════════════════════════════════════════════════
% BEHAVIORAL FEATURES
% ══════════════════════════════════════════════════════════════════
\node[inpMain, minimum width=12.4cm] (input) at (0, -0.4)
  {{\textcolor{neutD}{\faProjectDiagram}}\enspace {Instrumented Feature Vector}\enspace($H_{\text{attn}}$,\; $\cos(Q,K)$,\; $\|\nabla\|$,\; CKA,\; $\ldots$)};

\draw[mainArrow] (dynamic.south) -- (dynamic.south |- input.north);

% ══════════════════════════════════════════════════════════════════
% FPG ENCODING — NOVEL (terracotta accent, dashed)
% ══════════════════════════════════════════════════════════════════

% Row 1
\node[grp] (g1)  at (-5.0, -2.5) {Attention Stat.};
\node[grp] (g2)  at (-2.5, -2.5) {QKV Alignment};
\node[grp] (g3)  at ( 0.0, -2.5) {Score Stat.};
\node[grp] (g4)  at ( 2.5, -2.5) {Positional};
\node[grp] (g5)  at ( 5.0, -2.5) {FFN Output};
% Row 2
\node[grp] (g6)  at (-5.0, -3.3) {LayerNorm};
\node[grp] (g7)  at (-2.5, -3.3) {Residual Stream};
\node[grp] (g8)  at ( 0.0, -3.3) {Embedding};
\node[grp] (g9)  at ( 2.5, -3.3) {Repr.\ Drift};
\node[grp] (g10) at ( 5.0, -3.3) {Task Metrics};
% Row 3
\node[grp] (g11) at (-2.5, -4.1) {Training Dyn.};
\node[grp] (g12) at ( 0.0, -4.1) {Kernel Timing};
\node[grp] (g13) at ( 2.5, -4.1) {Cache Diag.};

% FPG directed edges (arrowheads show propagation direction)
% Horizontal: left → right (fault propagation across groups)
\foreach \a/\b in {g1/g2, g2/g3, g3/g4, g4/g5,
                   g6/g7, g7/g8, g8/g9, g9/g10} {
  \draw[fpgEdge] (\a) -- (\b);
}
% Vertical: top → bottom (fault propagation down layers)
\foreach \a/\b in {g1/g6, g2/g7, g3/g8, g4/g9, g5/g10,
                   g7/g11, g8/g12, g9/g13} {
  \draw[fpgEdge] (\a) -- (\b);
}

% Background box (terracotta dashed — novel)
\begin{scope}[on background layer]
  \coordinate (enc_top_pad) at ([yshift=18pt]g1.north);
  \node[draw=neutD!40, fill=neutL!50, rounded corners=5pt,
        inner xsep=12pt, inner ysep=7pt, densely dashed,
        fit=(enc_top_pad)(g1)(g5)(g6)(g10)(g11)(g13), line width=0.7pt] (enc) {};
\end{scope}

% Arrow: Input → Encoding
\draw[mainArrow] (input.south) -- (input.south |- enc.north)
  node[annot, midway, right=4pt] {structured by architectural dependencies};

% Encoding label (terracotta = novel)
\node[font=\sffamily\normalsize, text=neutD, fill=neutL!60,
      inner xsep=5pt, inner ysep=2pt]
  at ([yshift=-10pt]enc.north)
  {{\textcolor{neutD}{\faProjectDiagram}}\enspace Feature Encoding via the Fault Propagation Graph (FPG)};

% ══════════════════════════════════════════════════════════════════
% LEVEL 1 — DETECTION 
% ══════════════════════════════════════════════════════════════════
\node[stpbox, minimum width=8.0cm, densely dashed] (det) at (-0.5, -5.7)
  {{\textcolor{navyM}{\faSearch}}\enspace\textbf{Fault Detection} \enspace\textbar\enspace
   \textit{Is it faulty?}};

\draw[mainArrow] (det.north |- enc.south) -- (det.north);

% ── "No" branch → Healthy terminal (left side, away from losses) ──
\node[healthy] (healthyNode) at ([xshift=-70pt]det.west)
  {{\textcolor{sageD}{\faCheckCircle}}\enspace No Fault};
\draw[->, thick, arwC!60] ([xshift=-14pt]det.west) -- (healthyNode.east)
  node[annot, midway, above=2pt] {No};

% ══════════════════════════════════════════════════════════════════
% LEVEL 2 — CATEGORIZATION 
% ══════════════════════════════════════════════════════════════════

% Fault category chips (sage green)
\node[catchip] (h1)  at (-4.65, -8.0) {QKV Projection};
\node[catchip] (h2)  at (-1.55, -8.0) {Feed-Forward};
\node[catchip] (h3)  at ( 1.55, -8.0) {Score Computation};
\node[catchip] (h4)  at ( 4.65, -8.0) {Attention Masking};

\node[catchip] (h5)  at (-4.65, -8.8) {Layer Normalization};
\node[catchip] (h6)  at (-1.55, -8.8) {Positional Encoding};
\node[catchip] (h7)  at ( 1.55, -8.8) {Embedding Layer};
\node[catchip] (h8)  at ( 4.65, -8.8) {Output Projection};

\node[catchip] (h9)  at (-4.65, -9.6) {Residual Connection};
\node[catchip] (h10) at (-1.55, -9.6) {Kernel Runtime};
\node[catchip] (h11) at ( 1.55, -9.6) {Arch.\ Variant};
\node[catchip] (h12) at ( 4.65, -9.6) {KV Cache};

% Background (sage, dashed)
\begin{scope}[on background layer]
  \coordinate (cat_top_pad) at ([yshift=18pt]h1.north);
  \node[draw=navyM!50, fill=navyL!60, rounded corners=5pt,
        inner xsep=12pt, inner ysep=7pt, densely dashed,
        fit=(cat_top_pad)(h1)(h4)(h9)(h12), line width=0.7pt] (catbox_bg) {};
\end{scope}

% Label
\node[stlbl, fill=navyL!50, inner xsep=5pt, inner ysep=2pt]
  at ([yshift=-10pt]catbox_bg.north)
  {{\textcolor{navyM}{\faSitemap}}\enspace\textbf{Fault Categorization} \enspace\textbar\enspace
   {\normalfont\itshape What kind of fault is it?}};

\draw[condArrow] (det.south) -- (det.south |- catbox_bg.north)
  node[annot, midway, right=4pt] {Yes};

% ══════════════════════════════════════════════════════════════════
% LEVEL 3 — ROOT-CAUSE DIAGNOSIS
% ══════════════════════════════════════════════════════════════════
\node[stpbox, minimum width=9.5cm, densely dashed] (rc) at (0, -11.3)
  {{\textcolor{navyM}{\faCrosshairs}}\enspace\textbf{Root-Cause Diagnosis} \enspace\textbar\enspace
   \textit{What is the root cause within that category?}};

\draw[condArrow] (catbox_bg.south) -- (rc.north);

% ══════════════════════════════════════════════════════════════════
% TRAINING OBJECTIVES — right sidebar column
% ══════════════════════════════════════════════════════════════════

% Loss pill for Detection (from det.east, right side)
\node[losspill, anchor=west]
  (lossdet) at ([xshift=16pt]det.east)
  {{\textcolor{rustM}{\faCalculator}}\enspace$\mathcal{L}_{\text{det}}$\enspace binary CE};
\draw[densely dashed, rustD!40, line width=0.5pt, -]
  (det.east) -- (lossdet.west);

% Loss pill for Categorization (aligned to catbox)
\node[losspill, anchor=west]
  (losscat) at ([xshift=16pt]catbox_bg.east)
  {{\textcolor{rustM}{\faCalculator}}\enspace$\mathcal{L}_{\text{cat}}$\enspace multi-class CE};
\draw[densely dashed, rustD!40, line width=0.5pt, -]
  (catbox_bg.east) -- (losscat.west);

% Novel: Diagnostic Loss box
\node[lossbox, anchor=north west, inner xsep=8pt, inner ysep=6pt]
  at ([xshift=14pt, yshift=14pt]rc.north east)
  (lossboxN)
  {{\textcolor{rustD}{\faBalanceScale}}\enspace{\small\textcolor{rustD}{Root-Cause Separation Loss}}\\[6pt]
   $\mathcal{L}_{\text{diag}} = \mathcal{L}_{\text{rc}}
   + \beta\,\mathcal{L}_{\text{sc}}
   + \gamma\,\mathcal{L}_{\text{proto}}$\\[6pt]
   \begin{tabular}{@{}l@{\,}l@{}}
     {\textcolor{rustM}{\faAngleRight}}\,$\mathcal{L}_{\text{rc}}$ & root-cause CE\\[1pt]
     {\textcolor{rustM}{\faAngleRight}}\,$\mathcal{L}_{\text{sc}}$ & contrastive separation\\[1pt]
     {\textcolor{rustM}{\faAngleRight}}\,$\mathcal{L}_{\text{proto}}$ & prototypical alignment
   \end{tabular}};

% Dashed connector: Step 3 — Loss box (route around rc box)
\draw[densely dashed, rustD!40, line width=0.5pt, -]
  (rc.east) -- ++(0.3,0) |- (lossboxN.west);

% "Training Objectives" label (plain text, no box)
\node[font=\sffamily\small, text=rustD,
      inner xsep=2pt, inner ysep=2pt]
  (trainobj) at ([xshift=60pt, yshift=5pt]enc.south east)
  {{\textcolor{rustD}{\faCogs}}\enspace Training Objectives};

% Dashed connectors from Training Objectives to each loss element
% Route along a vertical spine on the right side of the loss pills, then branch left
% Single vertical spine at x = right edge of lossboxN + 8pt
\pgfmathsetmacro{\spineX}{8pt}
\coordinate (sp1) at ([xshift=8pt]lossboxN.north east |- lossdet.east);
\coordinate (sp2) at ([xshift=8pt]lossboxN.north east |- losscat.east);
\coordinate (sp3) at ([xshift=8pt]lossboxN.east);
% trainobj down to spine top
\draw[densely dashed, rustD!40, line width=0.5pt, -]
  (trainobj.south) |- (sp1);
% vertical spine: top to bottom
\draw[densely dashed, rustD!40, line width=0.5pt, -]
  (sp1) -- (sp3);
% horizontal branches into each loss
\draw[densely dashed, rustD!40, line width=0.5pt, -]
  (sp1) -- (lossdet.east);
\draw[densely dashed, rustD!40, line width=0.5pt, -]
  (sp2) -- (losscat.east);
\draw[densely dashed, rustD!40, line width=0.5pt, -]
  (sp3) -- (lossboxN.east);

% ══════════════════════════════════════════════════════════════════
% OUTPUT: EXPLANATION (solid, navy wash — conclusion emphasis)
% ══════════════════════════════════════════════════════════════════
\node[outbox, minimum height=2.0cm,
      inner xsep=6pt, inner ysep=6pt] (expl) at (0, -13.5) {%
  {{\textcolor{navyM}{\faLightbulb}}\enspace\textbf{Explanation via FPG}}\\[6pt]
  \begin{tabular}{@{}l@{\quad}l@{}}
  {Predicted root cause} \\ 
  {Most contributing features} 
  \end{tabular}};

% Arrow: Level 3 → Explanation
\draw[mainArrow] (rc.south) -- (expl.north); 

% ══════════════════════════════════════════════════════════════════
% Stage badges (navy — consistent)
% ══════════════════════════════════════════════════════════════════
\node[badge] at ([xshift=-14pt]det.west) {\footnotesize 1};
\node[badge] at ([xshift=-14pt]catbox_bg.west)   {\footnotesize 2};
\node[badge] at ([xshift=-14pt]rc.west)          {\footnotesize 3};

\end{tikzpicture}
\end{document}
